% Source: Romandini, 02_Introduction_Cryptography (+ lecture of 14 Nov 2026, morning)
\section{Introduction to Cryptography}
\label{sec:crypto-intro}

\subsection{Where cryptography fits in security}

Of the three CIA properties (\cref{sec:security-basics}), cryptography is the main technical tool for \textbf{confidentiality}: data stays unreadable for anyone but the intended parties. It also supports two further properties:
\begin{itemize}
    \item \textbf{Integrity}: the receiver can \emph{detect} that someone has tampered with the data while it travelled. Cryptography cannot stop an attacker from changing bits on the wire; it makes the change visible.
    \item \textbf{Authenticity}, i.e.\ integrity of the \emph{sender}: the receiver can be sure that a message signed \enquote{Mario Rossi} was really produced by Mario Rossi and not forged by someone else on his behalf.
\end{itemize}
Cryptography does \textbf{not} provide availability. Defences against DoS (resource placement, scaling) or ransomware (backups) are system-level measures, not cryptographic algorithms.

\paragraph{Threat model.} Throughout the cryptography lectures the \emph{endpoints are trusted} and the \emph{channel is untrusted}. The sender and receiver (by convention \textbf{Alice} and \textbf{Bob}) are assumed to run secure machines (no malware on their computers), while everything that travels between them can be read, recorded, modified, or injected by the adversary. This is the classical Dolev--Yao view \cite{dolev1983}. The Internet is the textbook untrusted infrastructure: the only \enquote{secure} alternative would be a private cable laid between the two parties and guarded along its whole length, and even then anyone who taps it can listen; on wireless links an antenna in a public square is enough.

The historic motivation of cryptography is therefore \emph{private communication over a public channel}, an old problem driven above all by military needs.

\subsection{A short history}

\begin{table}[H]
\centering
\caption{Milestones in the history of cryptography.}
\label{tab:crypto-history}
\begin{tabularx}{\textwidth}{@{}l X@{}}
\toprule
Period & Milestone \\
\midrule
Antiquity & Scytale (Sparta, transposition: letters rearranged) and Caesar cipher (substitution: letters shifted), all by hand. \\
9th century & Al-Kindi describes \emph{frequency analysis}, which breaks simple substitution ciphers. \\
16th--19th century & Polyalphabetic ciphers (Vigenère), considered unbreakable for centuries and broken in the 19th century (Babbage, Kasiski). \\
1920s--1940s & Electro-mechanical rotor machines (Enigma), broken at Bletchley Park. \\
1949 & Shannon's \emph{Communication Theory of Secrecy Systems} \cite{shannon1949} turns cryptography into a science. \\
1970s & DES (1977), the first public standard cipher, and public-key cryptography (Diffie--Hellman and Merkle, 1976; RSA, 1977). \\
2001 & AES becomes the symmetric standard. \\
2024 & NIST publishes the first post-quantum standards. \\
\bottomrule
\end{tabularx}
\end{table}

\subsection{The paradigm shift: from obscurity to Kerckhoffs}

\begin{definition}[Security through obscurity]
A design in which the protection relies on the \emph{algorithm itself} being secret: anyone who learns how the transformation works can read every message.
\end{definition}

Classical ciphers worked this way, like the secret alphabet children invent: it is secure only until someone figures out the alphabet. At that point \emph{every} message ever sent is readable, including old ones the adversary has recorded. In their own time these systems did work (try breaking a Caesar ciphertext without knowing anything about the method), but the approach does not scale.

\begin{definition}[Kerckhoffs' principle]
A cryptosystem must remain secure even if everything about it, except the key, is public \cite{kerckhoffs1883}. In Shannon's words: \enquote{the enemy knows the system} \cite{shannon1949}.
\end{definition}

The algorithm is public (published, often open source, with all its mathematics); the only secret is a \emph{parameter} of the algorithm, the \textbf{key}. Why is this better, even though it hands the attacker the full blueprint? The lecturer gave two reasons:
\begin{enumerate}
    \item \textbf{Public scrutiny.} A public algorithm is analysed by the whole community; independent researchers can find weaknesses and say \enquote{stop using it}. Users do not have to trust the designer's claim that it is secure: they can check. Secret designs are reviewed by no one and usually fail as soon as they are reverse-engineered.
    \item \textbf{Recovery after compromise.} Suppose an adversary who has been recording traffic eventually learns the secret. If the secret is a key, you generate a new one and immediately regain the same security guarantees. If the secret is the algorithm, you must invent a whole new algorithm. A key is short and easy to change; an algorithm is not.
\end{enumerate}

\begin{remark}
Keys get lost or leaked all the time; Kerckhoffs' principle is what makes this a routine event rather than a catastrophe. When designing or evaluating a cipher, assume the attacker knows \emph{everything} except the key.
\end{remark}

\subsection{Passive and active attacks}

\begin{figure}[H]
\centering
\begin{tikzpicture}[
    pt/.style={circle, draw, fill=black!15, minimum size=0.55cm, inner sep=0pt},
    panel/.style={draw, rounded corners, inner sep=6pt},
    lab/.style={font=\footnotesize, align=center},
    >=Stealth]
    % normal flow
    \node[pt] (a0) at (0,0) {};
    \node[pt] (b0) at (2.2,0) {};
    \draw[->, thick] (a0) -- (b0);
    \node[lab, below=0.75cm] at (1.1,0) {(a) Normal flow};
    % interception
    \begin{scope}[xshift=3.6cm]
        \node[pt] (a1) at (0,0) {};
        \node[pt] (b1) at (2.2,0) {};
        \node[pt, fill=red!35] (x1) at (1.1,-0.9) {};
        \draw[->, thick] (a1) -- (b1);
        \draw[->, thick, red!70!black] (1.1,0) -- (x1);
        \node[lab, below=0.15cm of x1] {(b) Interception\\\textcolor{themeblue}{passive: confidentiality}};
    \end{scope}
    % modification
    \begin{scope}[xshift=7.2cm]
        \node[pt] (a2) at (0,0) {};
        \node[pt] (b2) at (2.2,0) {};
        \node[pt, fill=red!35] (x2) at (1.1,-0.9) {};
        \draw[->, thick] (a2) -| (x2);
        \draw[->, thick, red!70!black] (x2) |- (b2);
        \node[lab, below=0.15cm of x2] {(c) Modification\\\textcolor{red!70!black}{active: integrity}};
    \end{scope}
    % fabrication
    \begin{scope}[xshift=10.8cm]
        \node[pt] (a3) at (0,0) {};
        \node[pt] (b3) at (2.2,0) {};
        \node[pt, fill=red!35] (x3) at (1.1,-0.9) {};
        \draw[->, thick, red!70!black] (x3) |- (b3);
        \node[lab, below=0.15cm of x3] {(d) Fabrication\\\textcolor{red!70!black}{active: authenticity}};
    \end{scope}
\end{tikzpicture}
\caption{Attacks on a communication between a source (left) and a destination (right); the shaded node is the attacker.}
\label{fig:passive-active}
\end{figure}

\begin{definition}[Passive and active attacks]
\leavevmode
\begin{itemize}
    \item A \textbf{passive attack} only observes the channel (eavesdropping: listening and recording). The message still reaches the destination unchanged and neither party knows it was read. Property at risk: \textbf{confidentiality} (interception).
    \item An \textbf{active attack} interferes with the communication. The attacker either \emph{modifies} a message in transit (attack on \textbf{integrity}) or \emph{fabricates} a message in someone else's name (attack on \textbf{authenticity}).
\end{itemize}
\end{definition}

\begin{example}
Sending a message over an open Wi-Fi network without encryption: the bits travel through the air and anyone in range can copy them. Copying does not consume the message, so the receiver gets it normally: a passive attack. A man-in-the-middle who changes the amount in a bank transfer, or who sends a letter signed with someone else's name, performs an active attack.
\end{example}

\paragraph{Countermeasures.} There are always two options:
\begin{itemize}
    \item \textbf{Channel access control}: control physically who can reach the channel (the guarded cable). Effective but expensive, and impossible on wireless or on the Internet.
    \item \textbf{Cryptographic protection}, cheaper and usable on any channel:
    \begin{itemize}
        \item against passive attacks, an \emph{unintelligible representation} of the data: without the key the transformed data cannot be read (encryption);
        \item against active attacks, \emph{integrity and origin attestation}: mechanisms that let the receiver verify that the content is unchanged and comes from the claimed sender (MACs and digital signatures, \cref{sec:base-mechanisms,sec:asymmetric}).
    \end{itemize}
\end{itemize}

\subsection{Terminology and notation}

\begin{definition}[Basic terms]
\leavevmode
\begin{itemize}
    \item \textbf{Plaintext} $m$: the original message or data, input to the encryption algorithm.
    \item \textbf{Encryption algorithm}: transforms the plaintext to make it unintelligible (private).
    \item \textbf{Key} $k$: second input of the encryption algorithm, a \emph{parameter} that determines the exact transformation applied to the plaintext.
    \item \textbf{Ciphertext} $c$: the transformed plaintext, output of the encryption algorithm.
    \item \textbf{Decryption algorithm}: the inverse of encryption; from the ciphertext and the appropriate key it reproduces the plaintext.
    \item \textbf{Cryptanalyst} $X$: the adversary who tries to break the scheme.
\end{itemize}
\end{definition}

\begin{figure}[H]
\centering
\begin{tikzpicture}[
    box/.style={draw, rounded corners, fill=themeblue!8, minimum height=0.9cm, align=center, font=\small},
    >=Stealth]
    \node[box] (A) {Alice\\$\Enc_k : m \mapsto c$};
    \node[box, right=3.4cm of A] (B) {Bob\\$\Dec_k : c \mapsto m$};
    \draw[->, thick] (A) -- node[above, font=\footnotesize] {$c$} node[below, font=\footnotesize\itshape] {public channel} (B);
    \node[box, fill=red!12, below=1.0cm] (X) at ($(A)!0.5!(B)$) {$X$ (cryptanalyst)};
    \draw[->, dashed, red!70!black] ($(A)!0.5!(B)$) ++(0,-0.3) -- (X);
\end{tikzpicture}
\caption{The basic communication scenario.}
\label{fig:comm-scenario}
\end{figure}

\paragraph{Notation.} The slides write $C_k(m)$ and $D_k(c)$; these notes use $\Enc_k(m)$ and $\Dec_k(c)$.
\begin{itemize}
    \item $\Enc_k(m) = c$: encryption of $m$ with key $k$.
    \item $\Dec_k(c) = m$: decryption of $c$ with key $k$.
    \item Correctness: $\Dec_k$ is the inverse of $\Enc_k$, i.e.\ $\Dec_k(\Enc_k(m)) = m$ for every $m$.
    \item If also $\Enc_k(\Dec_k(m)) = m$, the two functions are \textbf{commutative}. This property matters for public-key schemes such as RSA, where it is what makes digital signatures possible (\cref{sec:asymmetric}).
\end{itemize}

\subsection{Two families of cryptography}

\begin{definition}[Symmetric and asymmetric cryptography]
\leavevmode
\begin{itemize}
    \item \textbf{Symmetric} (secret-key, private-key) cryptography: the \emph{same} key is used to encrypt and to decrypt, so it must be shared by sender and receiver. Encryption and decryption are often the same (or very similar) function, differing only in the order of the inputs.
    \item \textbf{Asymmetric} (public-key) cryptography: \emph{different} keys are used to encrypt and to decrypt, and the encryption and decryption functions are different.
\end{itemize}
\end{definition}

\subsection{Cryptanalysis and attack models}

\begin{definition}[Cryptanalysis]
The science of breaking cryptographic schemes. Cryptography and cryptanalysis together form \emph{cryptology}.
\end{definition}

Motives range from curiosity to espionage and malice. The cryptanalyst may be \emph{passive} (eavesdrops only) or \emph{active} (injects or modifies messages). By Kerckhoffs' principle the cryptanalyst always knows the encryption and decryption algorithms, plus anything that can be learnt by observing past communications between A and B (and by any other means, e.g.\ social engineering or guessing likely plaintexts). The only assumption that holds is that the \emph{key} is secret: if the attacker also knew the key, no system could be built.

What the cryptanalyst knows \emph{in addition} to the algorithms defines the attack model (\cref{tab:attack-models}). Going down the table the attacker knows more, and a modern cipher must resist all of these models.

\begin{table}[H]
\centering
\caption{Attack models, from weakest to strongest attacker.}
\label{tab:attack-models}
\begin{tabularx}{\textwidth}{@{}l X@{}}
\toprule
Attack & What the cryptanalyst has (besides the algorithms) \\
\midrule
Brute force & Nothing but a ciphertext: tries every possible key. \\
Ciphertext only & A collection of ciphertexts $c_1, \dots, c_n$. \\
Known plaintext & Ciphertexts plus some pairs $(m_i, c_i)$ (e.g.\ standard headers, greetings, file formats). \\
Chosen plaintext & Pairs $(m_i, c_i)$ where the $m_i$ are \emph{chosen by the attacker}, e.g.\ by tricking the sender or a server into encrypting messages of the attacker's choice. \\
Chosen ciphertext & In addition, the decryptions of ciphertexts of the attacker's choice. \\
\bottomrule
\end{tabularx}
\end{table}

\begin{remark}[Brute force]
Brute force is always possible: knowing the algorithm and one ciphertext, the attacker enumerates all keys until the output \enquote{makes sense}. Strictly speaking it needs a way to recognise the right plaintext, which is why the slide's \enquote{nothing other than the algorithms} really means \enquote{nothing other than the algorithms and the ciphertext}. It cannot be prevented, only made \emph{infeasible}: keys are around 128--256 bits (symmetric) or 2048--4096 bits (RSA), and physical limits on time and computing power mean nobody can try $2^{128}$ keys. A large key space is \emph{necessary but not sufficient}: a cipher with a huge key space can still be broken by smarter attacks (e.g.\ frequency analysis against a monoalphabetic substitution, which has $26! \approx 2^{88}$ keys).
\end{remark}

\subsection{Perfect and computational secrecy}

Can we build a \enquote{perfect} cipher, i.e.\ one whose confidentiality is always guaranteed? Two notions must be distinguished.

\begin{definition}[Perfect and computational secrecy]
\leavevmode
\begin{itemize}
    \item \textbf{Perfect secrecy}: confidentiality is guaranteed against an adversary with \emph{unlimited} resources. Formally (Shannon), the ciphertext leaks nothing about the plaintext: $\Pr[M = m \mid C = c] = \Pr[M = m]$ for every $m$ and $c$.
    \item \textbf{Computational secrecy}: confidentiality is guaranteed only if the adversary's \emph{resources are limited} (computing power, time, memory, bandwidth, \dots). An adversary with realistic resources can do no better than guessing.
\end{itemize}
\end{definition}

Shannon (1949) proved that perfect secrecy is achievable, but only if the key is \textbf{truly random}, \textbf{at least as long as the message}, and \textbf{never reused}: the one-time pad (\cref{sec:symmetric}). With a key shorter than the message, perfect secrecy is impossible, because a brute-force adversary can try keys until the text makes sense.

\begin{remark}[What \enquote{perfect} means]
The lecturer stressed this point: perfect secrecy does \emph{not} mean that brute force is impossible. You can still try every key, but with a one-time pad every plaintext of the right length comes out of some key, so all candidates are equally plausible and the ciphertext reveals nothing about which one was sent. The one-time pad is impractical: the key must be exchanged securely, it must be as long as all the traffic, and it must be truly random rather than pseudo-random, which ordinary computers cannot easily produce.
\end{remark}

In practice we use ciphers with \textbf{computational secrecy}. If breaking a cipher takes a billion years of computation, or all the computing power in the world, it is \enquote{breakable} in theory but \emph{secure enough}. We want something usable, not something unbreakable.

\subsection{Symmetric cryptography}

\paragraph{Metaphor.} A chest sent over a public channel, closed with a padlock: Alice and Bob each hold an identical padlock and an identical key, and whoever has the key can both lock and unlock the chest. The catch is in the slide's caption: they \emph{need to agree on a common key before use}.

\paragraph{Features.}
\begin{itemize}
    \item Sender and receiver both know a common secret key $k$, both can encrypt and decrypt, and both promise to keep $k$ secret.
    \item The two parties must \textbf{agree on the key before use}, and they cannot do it over the public channel they want to protect (an eavesdropper would learn the key). Classically they use an \textbf{out-of-band} channel used only once: meeting in person, a courier, a sealed envelope. It is slow and expensive, but its cost is amortised over the many later exchanges on the public channel. Some processes still require people to physically meet to exchange keys, but over the Internet this is infeasible in general; key exchange is solved with public-key techniques (\cref{sec:asymmetric}).
    \item The \textbf{key space must be very large} to resist brute force (necessary but not sufficient, see above).
    \item \textbf{Advantage}: symmetric algorithms are \emph{very efficient}, much faster than asymmetric ones.
\end{itemize}

\paragraph{Ciphers in use.} DES (56-bit key) and Triple-DES are obsolete; Blowfish and IDEA are legacy; \textbf{AES} is the current standard (details in \cref{sec:symmetric}).

\paragraph{Number of keys.} With $n$ parties who all want to talk privately in pairs, every pair needs its own key. Each receiver has $n-1$ possible senders, and each key is shared by two people, so the total is
\[
    \binom{n}{2} = \frac{n(n-1)}{2} = O(n^2)
\]
keys, and each party must store $n-1$ of them. This, together with key distribution, is the main defect of symmetric cryptography.

\subsection{Asymmetric (public-key) cryptography}

Is it possible to exchange information confidentially \emph{without} first agreeing on a key? The answer came in 1976 with Diffie--Hellman (and Merkle), a turning point with huge theoretical and economic consequences: public-key cryptography is the foundation of e-commerce and of secure communication on the Internet.

\paragraph{Idea.} Break the symmetry between encrypting and decrypting: \emph{whoever knows how to encrypt must not know how to decrypt}.
\begin{itemize}
    \item The destination generates a key pair $(k_{\mathrm{priv}}, k_{\mathrm{pub}})$; the two keys are mathematically linked.
    \item $k_{\mathrm{priv}}$ is kept secret by the destination and used to \textbf{decrypt}.
    \item $k_{\mathrm{pub}}$ is published (it is a \enquote{key}, but not secret at all) and used by \emph{everyone} to \textbf{encrypt} messages for the destination.
    \item It must be computationally hard to derive $k_{\mathrm{priv}}$ from $k_{\mathrm{pub}}$.
\end{itemize}
If Alice encrypts with Bob's public key, only Bob, the sole holder of the matching private key, can decrypt.

\begin{definition}[One-way trapdoor function]
A function that is \emph{easy to compute}, \emph{hard to invert} in general, but \emph{easy to invert given some extra information} (the trapdoor, which plays the role of the private key).
\end{definition}

\begin{example}[Metaphors and the classic instance]
A padlock is easy to snap shut but hard to open without the key. A letterbox is easy to drop a letter into, but hard to take one out of unless you have the key. Mathematically: given two primes $p$ and $q$, computing $n = pq$ is easy; given only $n$, finding $p$ is hard, unless you already know $q$, in which case $p = n/q$. This is the basis of RSA (\cref{sec:asymmetric}).
\end{example}

\paragraph{Number of keys.} Each receiver has one public key, which all $n-1$ potential senders use. For $n$ parties, $n$ key pairs ($n$ public plus $n$ private keys) suffice: $O(n)$ instead of $O(n^2)$.

\begin{example}[Symmetric vs asymmetric key count]
For $n = 1000$ employees, pairwise symmetric keys: $1000 \cdot 999 / 2 = 499{,}500$ keys. Public-key: $1000$ key pairs.
\end{example}

\begin{table}[H]
\centering
\caption{Symmetric vs asymmetric cryptography.}
\label{tab:sym-vs-asym}
\begin{tabularx}{\textwidth}{@{}l X X@{}}
\toprule
 & Symmetric & Asymmetric \\
\midrule
Keys & One shared secret key & Key pair: public + private \\
Prior agreement & Needed (out-of-band) & Not needed \\
Keys for $n$ parties & $n(n-1)/2$ & $n$ pairs \\
Speed & Very fast & Much slower \\
Basis & Substitutions, permutations, XOR & Hard mathematical problems (trapdoors) \\
Examples & AES, ChaCha20 & RSA, ECC \\
\bottomrule
\end{tabularx}
\end{table}

\begin{remark}
In practice the two families are combined (\emph{hybrid encryption}): public-key cryptography is used to agree on or transport a fresh symmetric key, and the bulk data is then encrypted symmetrically for speed. TLS works this way.
\end{remark}

\subsection{Properties of modern cryptography}

Modern cryptography is much more than confidentiality. Three further properties are required:
\begin{itemize}
    \item \textbf{Integrity}: the receiver can determine whether the message has been tampered with (modified or replaced).
    \item \textbf{Authentication}: the receiver can ascertain that the message was sent by the presumed sender.
    \item \textbf{Non-repudiation}: the sender cannot later deny having sent the message, and the receiver can convince a \emph{third party} (a judge) that the sender really sent it. This is obtained with \textbf{digital signatures}.
\end{itemize}

\begin{remark}
Non-repudiation needs asymmetric cryptography. With a shared symmetric key (e.g.\ a MAC), either party could have produced the authentication tag, so the receiver cannot prove to a judge that the sender, and not the receiver, created the message. Only the holder of a private key can produce a digital signature.
\end{remark}

\subsection{The quantum threat: harvest now, decrypt later}

\begin{itemize}
    \item A large quantum computer running \textbf{Shor's algorithm} (1994) could factor integers and compute discrete logarithms efficiently, breaking today's public-key cryptography (RSA, Diffie--Hellman, ECC).
    \item \textbf{Symmetric ciphers} are much less affected: Grover's algorithm only gives a quadratic speed-up on key search, so \emph{doubling the key length} (e.g.\ AES-128 $\to$ AES-256) restores the security level.
    \item \textbf{Harvest now, decrypt later}: adversaries are \emph{already} recording encrypted traffic, to decrypt it once a powerful enough quantum computer exists. The problem is not \enquote{in 20 years we will need to adapt}: data that must stay confidential for many years is at risk today.
    \item No such machine exists today. The answer is \textbf{post-quantum cryptography} (first NIST standards in 2024), covered in \cref{sec:asymmetric}.
\end{itemize}

\begin{remark}[AI and cryptanalysis]
The lecturer also noted that AI is being applied to the mathematical problems behind public-key cryptography (e.g.\ factoring). To be precise, what would break RSA is an efficient \emph{factoring} algorithm, not a way of generating primes: generating large primes is already easy, and is exactly how RSA keys are made. A proof that P\,=\,NP with a practical algorithm would also break it. No such result exists today.
\end{remark}

\begin{summary}
\begin{itemize}
    \item Cryptography provides confidentiality, integrity (tamper \emph{detection}) and authenticity, never availability. Threat model: trusted endpoints, untrusted channel.
    \item Kerckhoffs: everything is public except the key, because public algorithms get reviewed and a leaked key is easy to replace.
    \item Passive attacks (interception) threaten confidentiality; active attacks threaten integrity (modification) and authenticity (fabrication).
    \item Attack models: brute force, ciphertext only, known plaintext, chosen plaintext, chosen ciphertext.
    \item Perfect secrecy (one-time pad) is possible but impractical; real ciphers offer computational secrecy.
    \item Symmetric: one shared key, fast, but needs prior key agreement and $n(n-1)/2$ keys. Asymmetric: key pair based on a one-way trapdoor function, $n$ pairs, no prior agreement.
    \item Modern goals: integrity, authentication, non-repudiation (digital signatures).
    \item Shor breaks public-key crypto; symmetric crypto survives by doubling key lengths. Harvest now, decrypt later makes it urgent.
\end{itemize}
\end{summary}

\subsection*{Exercises}

\begin{exercise}
State Kerckhoffs' principle and give two reasons why it is preferable to security through obscurity.
\end{exercise}
\begin{solution}
A cryptosystem must remain secure even if everything about it except the key is public. Reasons: (1) a public algorithm can be analysed by the whole community, so weaknesses are found and reported, and users need not trust the designer; (2) if the secret leaks, a key can be replaced immediately and the same guarantees are restored, whereas a leaked secret algorithm requires designing a new algorithm, and all past messages are exposed.
\end{solution}

\begin{exercise}
For each scenario, say whether the attack is passive or active and which property it targets: (a) sniffing unencrypted traffic on a café Wi-Fi; (b) altering the IBAN in a payment order in transit; (c) sending an e-mail that pretends to come from the CEO.
\end{exercise}
\begin{solution}
(a) Passive, confidentiality (interception). (b) Active, integrity (modification). (c) Active, authenticity (fabrication).
\end{solution}

\begin{exercise}
An attacker convinces a web server to encrypt, with its secret key, messages chosen by the attacker, and observes the resulting ciphertexts. Which attack model is this? Which model applies if instead the attacker only knows that every encrypted message starts with a fixed protocol header?
\end{exercise}
\begin{solution}
The first is a chosen-plaintext attack (the attacker selects the $m_i$ and obtains the $c_i$). The second is a known-plaintext attack: the attacker knows part of the plaintext for each ciphertext but did not choose it.
\end{solution}

\begin{exercise}
A company of 200 employees wants every pair of employees to communicate confidentially. How many keys are needed with symmetric cryptography, and how many with public-key cryptography? What other problem does the symmetric solution have?
\end{exercise}
\begin{solution}
Symmetric: $200 \cdot 199 / 2 = 19{,}900$ keys (each employee stores 199). Public-key: 200 key pairs. Moreover, each symmetric key must be agreed on in advance over a secure out-of-band channel, which does not scale.
\end{solution}

\begin{exercise}
Distinguish perfect secrecy from computational secrecy. Why is a cipher with a 128-bit key used for a 1 MB file not perfectly secret, and why do we use it anyway?
\end{exercise}
\begin{solution}
Perfect secrecy means the ciphertext reveals nothing about the plaintext even to an unbounded adversary; computational secrecy holds only against adversaries with limited resources. With a key much shorter than the message, an unbounded adversary could try all $2^{128}$ keys and recognise the one that yields a meaningful file, so perfect secrecy is impossible (Shannon). We use it because trying $2^{128}$ keys is far beyond any realistic resources, and a perfectly secret one-time pad would need a truly random, never-reused 1 MB key exchanged in advance.
\end{solution}

\begin{exercise}
Explain \enquote{harvest now, decrypt later}. Why does it make the migration to post-quantum cryptography urgent even though no cryptographically relevant quantum computer exists yet, and why are symmetric ciphers less concerned?
\end{exercise}
\begin{solution}
Adversaries record encrypted traffic today and store it until a quantum computer running Shor's algorithm can break the public-key cryptography (RSA, ECC) that protected the key exchange; then they decrypt the stored data. Any data that must remain confidential longer than the time until such machines exist is therefore already at risk. Symmetric ciphers are only weakened by Grover's quadratic speed-up, so doubling the key length (AES-256) is enough.
\end{solution}
